% LaTeX Chapter Template for Double ML Book
%
% Quick reference guide for writing chapters
%
% USAGE:
%   1. Copy this file: cp chapter_template.tex chapter_NN.tex
%   2. Replace NN with chapter number
%   3. Fill in content following examples below
%   4. Add to main.tex: \include{chapters/chapter_NN}
%   5. Compile: make

\chapter{Chapter Title Here}

% ============================================================================
% INTRODUCTION
% ============================================================================

\section{Introduction}

Start with motivation and overview of the chapter. Explain:
\begin{itemize}
    \item Why this topic matters
    \item What you'll learn
    \item How it connects to previous/future chapters
\end{itemize}

% ============================================================================
% MAIN CONTENT SECTIONS
% ============================================================================

\section{Main Section Title}

Your prose content here with inline math like $\tau = \E[Y_i(1) - Y_i(0)]$.

\subsection{Subsection Title}

More detailed content.

% ----------------------------------------------------------------------------
% DEFINITIONS
% ----------------------------------------------------------------------------

\begin{definition}[Name of Definition]
\label{def:short_name}
A \textbf{definition} provides precise terminology.

For example, the Average Treatment Effect (ATE) is:
\[
\tau = \E[Y_i(1) - Y_i(0)]
\]
where $Y_i(1)$ is the potential outcome under treatment and $Y_i(0)$ under control.
\end{definition}

% Reference the definition with: see Definition~\ref{def:short_name}

% ----------------------------------------------------------------------------
% THEOREMS
% ----------------------------------------------------------------------------

\begin{theorem}[Main Result]
\label{thm:main_result}
If assumptions A1-A3 hold, then:
\[
\sqrt{n}(\hat{\tau} - \tau_0) \convd N(0, \Sigma)
\]
\end{theorem}

\begin{proof}
Proof structure:
\begin{enumerate}
    \item \textbf{Step 1}: Establish preliminary result
    \[
    \text{Some equation}
    \]

    \item \textbf{Step 2}: Main argument

    Using the law of iterated expectations:
    \begin{align}
        \E[\psi(W_i; \tau_0, \eta_0)]
        &= \E\left[\E[\psi(W_i; \tau_0, \eta_0) \mid X_i]\right] \\
        &= 0
    \end{align}

    \item \textbf{Step 3}: Conclusion

    Therefore, the result holds.
\end{enumerate}
\end{proof}

% ----------------------------------------------------------------------------
% EXAMPLES
% ----------------------------------------------------------------------------

\begin{example}[Motivating Example]
\label{ex:motivation}
Consider an insurance pricing scenario where:
\begin{itemize}
    \item $Y_i$: customer churn (0/1)
    \item $T_i$: competitor price increase (\$)
    \item $X_i$: customer age, income, coverage
\end{itemize}

We want to estimate the causal effect of $T_i$ on $Y_i$.
\end{example}

% ----------------------------------------------------------------------------
% PYTHON CODE
% ----------------------------------------------------------------------------

\subsection{Python Implementation}

Here's how to implement this in Python with EconML:

\begin{minted}{python}
import numpy as np
from econml.dml import LinearDML
from sklearn.ensemble import RandomForestRegressor

# Set seed for reproducibility
np.random.seed(42)
n = 1000

# Generate data
X = np.random.randn(n, 5)
T = X[:, 0] + X[:, 1] + np.random.randn(n)
Y = 2.5 * T + X[:, 0]**2 + np.random.randn(n)

# Estimate DML
dml = LinearDML(
    model_y=RandomForestRegressor(n_estimators=100, random_state=42),
    model_t=RandomForestRegressor(n_estimators=100, random_state=42),
    cv=5,
    random_state=42
)

dml.fit(Y, T, X=X, W=None)

# Get treatment effect
tau_hat = dml.effect(X).mean()
print(f"Estimated ATE: {tau_hat:.3f}")
\end{minted}

% ----------------------------------------------------------------------------
% REMARKS AND NOTES
% ----------------------------------------------------------------------------

\begin{remark}
Important insights or caveats can go in remark blocks.

For example: The assumption of unconfoundedness cannot be tested with data alone.
\end{remark}

% ============================================================================
% EXERCISES
% ============================================================================

\section{Exercises}

\begin{exercise}[Conceptual Understanding]
\label{ex:conceptual}
Explain why we need both unconfoundedness and overlap for ATE identification.
\end{exercise}

\textbf{Solution to Exercise~\ref{ex:conceptual}:}

\begin{enumerate}
    \item \textbf{Unconfoundedness}: Ensures no omitted confounders

    Without this, we cannot identify the causal effect.

    \item \textbf{Overlap}: Ensures we have data for both treated and control groups across all covariate values

    Without this, we'd be extrapolating.
\end{enumerate}

\begin{exercise}[Mathematical Derivation]
\label{ex:derivation}
Prove that under unconfoundedness and overlap:
\[
\tau = \E\left[\frac{T_i Y_i}{e(X_i)} - \frac{(1-T_i)Y_i}{1-e(X_i)}\right]
\]
where $e(x) = \Prob(T_i=1 \mid X_i=x)$.
\end{exercise}

\textbf{Solution to Exercise~\ref{ex:derivation}:}

[Detailed mathematical proof here]

% ============================================================================
% SUMMARY
% ============================================================================

\section{Summary}

Recap the key points:
\begin{enumerate}
    \item First main insight
    \item Second main insight
    \item Third main insight
\end{enumerate}

\textbf{Key takeaways:}
\begin{itemize}
    \item Practical point 1
    \item Practical point 2
    \item Practical point 3
\end{itemize}

% ============================================================================
% ROADMAP TO NEXT CHAPTER
% ============================================================================

\section{Roadmap to Chapter [N+1]}

Preview the next chapter:

Chapter [N+1] extends these ideas to [next topic]. We will develop:
\begin{itemize}
    \item New concept 1
    \item New concept 2
    \item Application to [domain]
\end{itemize}

% ============================================================================
% QUICK REFERENCE GUIDE
% ============================================================================

% COMMON MATH COMMANDS:
%   \E[X]           Expectation
%   \Var(X)         Variance
%   \Cov(X,Y)       Covariance
%   \Prob(A)        Probability
%   \R              Real numbers
%   \convd          Convergence in distribution
%   \convp          Convergence in probability
%   \indicator{A}   Indicator function
%   \norm{x}        Norm
%
% EQUATION ENVIRONMENTS:
%   \[ ... \]                 Unnumbered equation
%   \begin{equation} ... \end{equation}        Numbered equation
%   \begin{align} ... \end{align}              Numbered multi-line
%   \begin{align*} ... \end{align*}            Unnumbered multi-line
%
% CROSS-REFERENCES:
%   \label{eq:name}          Label an equation
%   \ref{eq:name}            Reference equation number
%   \eqref{eq:name}          Reference with parentheses: (1.2)
%
% CITATIONS:
%   \cite{author2024}        Basic citation
%   \cite[Theorem 1]{author2024}   Citation with note
%
% THEOREM ENVIRONMENTS:
%   definition, theorem, lemma, proposition, corollary
%   example, exercise, remark, note
%
% CODE BLOCKS:
%   \begin{minted}{python} ... \end{minted}
%   \begin{minted}{r} ... \end{minted}
%   \begin{minted}{julia} ... \end{minted}
